\section{Conclusions}
\label{sec:conclusions}

This paper introduced a structured methodology for AI-assisted
conceptual aircraft design in which a large-language-model agent works
as an engineering collaborator inside a governed, version-controlled
design repository. The methodology rests on a strict separation of
roles --- deterministic, configured, regression-tested physics code as
the sole source of engineering numbers; the AI as the implementer and
propagator of design changes; humans as the source of requirements,
directives, review, and acceptance --- enforced by four transferable
mechanisms: a machine-readable project constitution, a one-way
analysis pipeline with typed hand-offs, pinned design-point regression
tests with full-pipeline continuous integration, and architecture
decision records with git-level provenance of AI contribution.

Applied to the conceptual design of a small electric tail-sitter VTOL
UAV, the methodology carried the program from mission requirements to
a converged, commercially frozen 2.52\,kg design point --- sixteen
analysis stages spanning sizing, aerodynamics, control authority,
vibration, thermal, CAD, component selection, and higher-fidelity
hand-off --- in three calendar weeks of part-time work, with thirteen
fully consistent tagged design snapshots and eighteen recorded
architecture decisions. Process evidence extracted from the git
history shows 64\% of design-phase commits co-authored by the AI, all
merged through human review, and documents which safety layer caught
each class of fault: automated pins caught silent drift, versioned
hand-off schemas caught cross-tool integration errors, and independent
expert review caught the conceptual errors no automated check could
see.

The answer to the research question is therefore affirmative but
conditional: AI can accelerate conceptual aircraft design by roughly
an order of magnitude in delivered scope per unit time \emph{when} it
is deployed at the workflow level under structural guardrails, rather
than as a source of engineering answers. The acceleration was used not
to cut corners but to make honesty cheap --- model refinements that
grew the design mass were adopted within days, and unfavorable margins
are carried openly as standing findings with named retirement paths.

Future work proceeds along the hand-offs the methodology already
exposes: closing the external vortex-lattice/blade-element and CFD
loops on the frozen geometry, PX4-based flight validation of the
hover/transition control architecture, folding measured component
properties (notably pack internal resistance) back into the closure
after procurement, and applying the framework to a second, dissimilar
airframe to test aircraft-independence. On the methodological side,
controlled multi-team studies are needed to separate the contributions
of the AI agent, the governance structure, and team expertise to the
observed acceleration.
